\documentclass[10pt,aspectratio=43,mathserif]{beamer}

\usepackage{seu}
\usepackage{xeCJK}
\usepackage{amsmath,amsfonts,amssymb,bm}
\usepackage{color}
\usepackage{graphicx,hyperref,url}
\usepackage{etoolbox}

\setsansfont[Path=fonts/]{Helvetica}
\beamertemplateballitem
\setbeamertemplate{footline}{}

% 正文统一使用较小字号；标题页、章节标题与 frame 标题保持主题默认字号
\setbeamerfont{normal text}{size=\footnotesize}
\setbeamerfont{block title}{size=\footnotesize}
\setbeamerfont{block title alerted}{size=\footnotesize}
\setbeamerfont{block body}{size=\footnotesize}
\setbeamerfont{block body alerted}{size=\footnotesize}
\setbeamerfont{itemize/enumerate body}{size=\footnotesize}
\setbeamerfont{itemize/enumerate subbody}{size=\scriptsize}

% 展示公式统一缩小一级，行内公式保持正文大小
\AtBeginEnvironment{displaymath}{\footnotesize}
\AtBeginEnvironment{equation}{\footnotesize}
\AtBeginEnvironment{equation*}{\footnotesize}
\AtBeginEnvironment{align}{\footnotesize}
\AtBeginEnvironment{align*}{\footnotesize}

% 将 block 改为无边框的“标题 + 正文”结构
\setbeamertemplate{block begin}{%
    \par\vspace{0.35em}%
    {\usebeamerfont{block title}\usebeamercolor[fg]{structure}\insertblocktitle\par}%
    \vspace{0.15em}%
    \usebeamerfont{block body}\usebeamercolor[fg]{normal text}%
}
\setbeamertemplate{block end}{\par\vspace{0.35em}}
\setbeamertemplate{block alerted begin}{%
    \par\vspace{0.35em}%
    {\usebeamerfont{block title alerted}\usebeamercolor[fg]{alerted text}\insertblocktitle\par}%
    \vspace{0.15em}%
    \usebeamerfont{block body alerted}\usebeamercolor[fg]{normal text}%
}
\setbeamertemplate{block alerted end}{\par\vspace{0.35em}}

% 图片暂未插入时使用的统一占位提示
\newcommand{\imageplaceholder}[2][2.5cm]{%
    \parbox[c][#1][c]{\linewidth}{%
        \centering\scriptsize\textcolor{black!60}{\textbf{图片占位：#2}}%
    }%
}

\AtBeginSection[]
{
    \begin{frame}<beamer>
        \frametitle{\textbf{目录}}
        \textbf{\tableofcontents[currentsection]}
    \end{frame}
}

\title[短标题]{\fontsize{13pt}{18pt}\selectfont 完整报告标题}
\author[页脚短作者]{作者甲, 作者乙 \\ \medskip}
\institute[短机构名]{完整机构名}
\date[\today]{\today}

\begin{document}

\begin{frame}
    \titlepage
\end{frame}

\section*{目录}
\begin{frame}
    \frametitle{\textbf{目录}}
    \textbf{\tableofcontents}
\end{frame}

\section[第一部分]{第一部分标题}

\begin{frame}
    \frametitle{\textbf{页面标题}}
    \begin{columns}[T]
        \column{.46\textwidth}
        \begin{block}{\textbf{块标题}}
            \begin{itemize}
                \item 要点一
                \item 要点二
            \end{itemize}
        \end{block}

        \column{.50\textwidth}
        \begin{block}{\textbf{块标题}}
            展示公式示例：
            \[
                V(s)=\mathbb{E}\left[\sum_{t=0}^{\infty}\gamma^t r_t\right]
            \]
        \end{block}
    \end{columns}
\end{frame}

\begin{frame}
    \frametitle{\textbf{图片占位示例}}
    \begin{columns}[T]
        \column{.48\textwidth}
        \begin{block}{\textbf{左侧说明}}
            \begin{itemize}
                \item 图文双栏时使用该版式
                \item 图片统一存放在 \texttt{figures/} 下
            \end{itemize}
        \end{block}

        \column{.48\textwidth}
        \centering
        \imageplaceholder{figures/example.png}
    \end{columns}
\end{frame}

\section[第二部分]{第二部分标题}

\begin{frame}
    \frametitle{\textbf{列表与要点}}
    \begin{block}{\textbf{有明确顺序时}}
        \begin{enumerate}
            \item 第一步
            \item 第二步
            \item 第三步
        \end{enumerate}
    \end{block}
\end{frame}

\section*{}
\begin{frame}
    \centering
    \LARGE\textbf{\quad Thanks for Listening.}
\end{frame}

\end{document}
